\section{Related Work}
\label{sec:related}

\paragraph{Video retrieval augmented generation.}
RAG grounds generation in retrieved evidence. For long videos, Video-RAG~\cite{luo2024videorag} augments a VLM with
visually aligned auxiliary texts, such as speech transcripts, on-screen text, and detected objects, retrieved for each
query. VideoRAG~\cite{ren2025videorag} instead builds a graph-based knowledge index across videos, and
Vgent~\cite{shen2025vgent} represents a video as a graph of clips that preserves their semantic and temporal
dependencies, followed by a structured reasoning step that filters retrieval noise. These systems, however, focus on
what to index and retrieve, and pass every retrieved item to the VLM at its full cost. \name instead decides how many
tokens each retrieved item receives, and can be placed after any of these retrievers. \ayush{add more cites, maybe ~6 more, mention AVA}

\paragraph{Keyframe selection.}
Most VLMs sample frames uniformly and therefore often miss the content that a query needs. Query-aware methods instead
select frames by jointly optimizing their relevance and temporal coverage~\cite{tang2025aks,peng2026qca}, by greedy
marginal utility~\cite{shi2026markey}, by learning to query informative combinations of
frames~\cite{yu2024framevoyager}, or with the attention of a small VLM used as a selector~\cite{wang2026dafs}. All of
these methods, however, return a subset of frames, each at the same token cost. \name instead treats the token count
of each frame as a decision variable.

